\section{Uncertainty Relations}

\SourcePage{67}{70}
Let us derive the commutation rules for the momentum and coordinate operators.
Successive differentiation with respect to one of the variables $x,y,z$ and
multiplication by another of them gives a result independent of the order of
these operations. Hence
\begin{equation}
  \hat p_xy-y\hat p_x=0,\qquad \hat p_xz-z\hat p_x=0,
  \tag{16.1}
\end{equation}
and similarly for $\hat p_y$ and $\hat p_z$.

To derive the commutation rule for $\hat p_x$ and $x$, write
\[
  (\hat p_xx-x\hat p_x)\psi
  =-i\hbar\frac{\partial}{\partial x}(x\psi)
  +i\hbar x\frac{\partial\psi}{\partial x}=-i\hbar\psi.
\]
We see that the action of $\hat p_xx-x\hat p_x$ reduces to multiplication of
the function by $-i\hbar$; the same is of course true for the commutation of
$\hat p_y$ with $y$ and of $\hat p_z$ with $z$. Thus\footnote{These
relations, discovered in matrix form by \emph{Heisenberg} in 1925, served as
the starting point in the creation of quantum mechanics.}
\begin{equation}
  \hat p_xx-x\hat p_x=-i\hbar,\qquad
  \hat p_yy-y\hat p_y=-i\hbar,\qquad
  \hat p_zz-z\hat p_z=-i\hbar.
  \tag{16.2}
\end{equation}
Relations (16.1) and (16.2) can be written together as
\begin{equation}
  \hat p_i x_k-x_k\hat p_i=-i\hbar\delta_{ik},
  \qquad i,k=x,y,z.
  \tag{16.3}
\end{equation}

Before examining the physical meaning of these relations and their
consequences, we give two formulas useful below. If $f(\mathbf r)$ is a
function of the coordinates, then
\begin{equation}
  \hat{\mathbf p}f(\mathbf r)-f(\mathbf r)\hat{\mathbf p}
  =-i\hbar\Grad{f}.
  \tag{16.4}
\end{equation}
Indeed,
\[
  (\hat{\mathbf p}f-f\hat{\mathbf p})\psi
  =-i\hbar[\Grad{(f\psi)}-f\Grad{\psi}]
  =-i\hbar\psi\Grad{f}.
\]
An analogous relation holds for the commutator with a function of the momentum
operator:
\begin{equation}
  f(\mathbf p)\hat{\mathbf r}-\hat{\mathbf r}f(\mathbf p)
  =-i\hbar\frac{\partial f}{\partial\mathbf p}.
  \tag{16.5}
\end{equation}
It may be derived in the same way as (16.4) by working in the momentum
representation and using (15.12) for the coordinate operators.

Relations (16.1) and (16.2) show that a particle coordinate along one axis can
have a definite value simultaneously with the momentum components along the
other two axes; a coordinate and the momentum component along that same axis%
\SourcePage{68}{71}
do not exist simultaneously. In particular, a particle cannot occupy a
definite point in space while also having a definite momentum $\mathbf p$.

Let the particle occupy a finite spatial region whose dimensions along the
three axes are of order $\Delta x$, $\Delta y$, $\Delta z$. Let its mean
momentum be $\mathbf p_0$. Mathematically, this means that the wave function
has the form
$\psi=u(\mathbf r)e^{i\mathbf p_0\cdot\mathbf r/\hbar}$, where
$u(\mathbf r)$ is appreciably different from zero only within the specified
region.

Expand $\psi$ in eigenfunctions of the momentum operator, that is, as a
Fourier integral. The coefficients $a(\mathbf p)$ in this expansion are given
by (15.10) as integrals of functions of the form
$u(\mathbf r)e^{i(\mathbf p_0-\mathbf p)\cdot\mathbf r/\hbar}$. For such an
integral to be appreciably different from zero, the periods of the oscillating
factor $e^{i(\mathbf p_0-\mathbf p)\cdot\mathbf r/\hbar}$ must not be small
compared with the dimensions $\Delta x$, $\Delta y$, and $\Delta z$ of the
region where $u(\mathbf r)$ differs from zero. Thus $a(\mathbf p)$ is
appreciably different from zero only for values of $\mathbf p$ satisfying
$(p_{0x}-p_x)\Delta x/\hbar\lesssim1,\ldots$ Since $|a(\mathbf p)|^2$
determines the probabilities of the various momentum values, the intervals of
$p_x,p_y,p_z$ in which $a(\mathbf p)$ is nonzero are precisely the intervals
in which the particle's momentum components may lie in the state considered.
Denoting these intervals by $\Delta p_x$, $\Delta p_y$, and $\Delta p_z$, we
therefore have
\begin{equation}
  \Delta p_x\Delta x\sim\hbar,\qquad
  \Delta p_y\Delta y\sim\hbar,\qquad
  \Delta p_z\Delta z\sim\hbar.
  \tag{16.6}
\end{equation}
These \emph{uncertainty relations} were established by \emph{Heisenberg} in
1927.

We see that the more accurately a particle coordinate is known (the smaller
$\Delta x$ is), the greater is the uncertainty $\Delta p_x$ in the momentum
component along the same axis, and conversely. In particular, if a particle
occupies a strictly definite point of space
($\Delta x=\Delta y=\Delta z=0$), then
$\Delta p_x=\Delta p_y=\Delta p_z=\infty$. All momentum values are then
equally probable. Conversely, if the particle has a strictly definite
momentum $\mathbf p$, all its positions in space are equally probable (this
also follows directly from the wave function (15.8), whose squared modulus is
entirely independent of the coordinates).

If the coordinate and momentum uncertainties are characterized by their
root-mean-square fluctuations,
\[
  \delta x=\sqrt{\overline{(x-\bar x)^2}},\qquad
  \delta p_x=\sqrt{\overline{(p_x-\bar p_x)^2}},
\]
\SourcePage{69}{72}
an exact bound can be given for the least possible value of their product
(\emph{H.~Weyl}).

Consider the one-dimensional case, a packet with wave function $\psi(x)$
depending on only one coordinate; for simplicity, suppose that the mean values
of $x$ and $p_x$ in this state are zero. Begin with the evident inequality
\[
  \int_{-\infty}^{\infty}
  \left|\alpha x\psi+\frac{d\psi}{dx}\right|^2dx\geq0,
\]
where $\alpha$ is an arbitrary real constant. In evaluating this integral,
note that
\[
  \int x^2|\psi|^2\,dx=(\delta x)^2,
\]
\[
  \int\left(x\frac{d\psi^*}{dx}\psi
  +x\psi^*\frac{d\psi}{dx}\right)dx
  =\int x\frac{d|\psi|^2}{dx}\,dx
  =-\int|\psi|^2\,dx=-1,
\]
\[
  \int\frac{d\psi^*}{dx}\frac{d\psi}{dx}\,dx
  =-\int\psi^*\frac{d^2\psi}{dx^2}\,dx
  =\frac{1}{\hbar^2}\int\psi^*\hat p_x^2\psi\,dx
  =\frac{1}{\hbar^2}(\delta p_x)^2,
\]
and obtain
\[
  \alpha^2(\delta x)^2-\alpha
  +\frac{(\delta p_x)^2}{\hbar^2}\geq0.
\]
For this quadratic trinomial in $\alpha$ to be nonnegative for every value of
$\alpha$, its discriminant must be negative. It follows that
\begin{equation}
  \delta x\,\delta p_x\geq\hbar/2.
  \tag{16.7}
\end{equation}
The least possible value of the product is $\hbar/2$.

This value is attained by wave packets described by functions of the form
\begin{equation}
  \psi=\frac{1}{(2\pi)^{1/4}\sqrt{\delta x}}
  \exp\left(\frac{i}{\hbar}p_0x-\frac{x^2}{4(\delta x)^2}\right),
  \tag{16.8}
\end{equation}
where $p_0$ and $\delta x$ are constants. The probabilities of the various
coordinate values in such a state are
\[
  |\psi|^2=\frac{1}{\sqrt{2\pi}\,\delta x}
  \exp\left(-\frac{x^2}{2(\delta x)^2}\right),
\]
that is, they have a Gaussian distribution about the origin (mean value
$\bar x=0$), with root-mean-square fluctuation $\delta x$. The wave function
in the momentum representation is
\[
  a(p_x)=\int_{-\infty}^{\infty}\psi(x)
  \exp\left(-i\frac{p_xx}{\hbar}\right)dx.
\]
\SourcePage{70}{73}
Evaluation of the integral gives an expression of the form
\[
  a(p_x)=\mathrm{const}\cdot
  \exp\left(-\frac{(\delta x)^2(p_x-p_0)^2}{\hbar^2}\right).
\]
The probability distribution of the momentum values, $|a(p_x)|^2$, is also
Gaussian, about the mean $\bar p_x=p_0$, and has root-mean-square fluctuation
$\delta p_x=\hbar/(2\delta x)$. Thus the product
$\delta p_x\delta x$ is exactly $\hbar/2$.

Finally, let us derive one more useful relation. Let $f$ and $g$ be two
physical quantities whose operators obey the commutation rule
\begin{equation}
  \hat f\hat g-\hat g\hat f=-i\hbar\hat c,
  \tag{16.9}
\end{equation}
where $\hat c$ is the operator of a physical quantity $c$. The factor $\hbar$
has been introduced on the right because in the classical limit
($\hbar\to0$) all operators of physical quantities reduce to multiplication
by those quantities and commute with one another. Thus, in the
``quasiclassical'' case, the right-hand side of (16.9) may be taken as zero in
the first approximation. In the next approximation, $\hat c$ may be replaced
by the operator of simple multiplication by $c$. We then have
\[
  \hat f\hat g-\hat g\hat f=-i\hbar c.
\]
This equation is exactly analogous to
$\hat p_xx-x\hat p_x=-i\hbar$, except that the quantity $\hbar c$ occurs in
place of the constant $\hbar$\footnote{The classical quantity $c$ is the
Poisson bracket of $f$ and $g$ (see the note on p.~46).}. By analogy with
$\Delta x\Delta p_x\sim\hbar$, we therefore conclude that in the
quasiclassical case the quantities $f$ and $g$ obey the uncertainty relation
\begin{equation}
  \Delta f\Delta g\sim\hbar c.
  \tag{16.10}
\end{equation}
In particular, if one of the quantities is the energy ($f\equiv H$) and the
operator of the other ($g$) has no explicit time dependence, then by (9.2)
$c=\dot g$, and the quasiclassical uncertainty relation is
\begin{equation}
  \Delta E\Delta g\sim\hbar\dot g.
  \tag{16.11}
\end{equation}
